\documentclass[a4paper]{ctexart}
\usepackage[margin=1.35cm,top=1.2cm,bottom=1.2cm]{geometry}
\usepackage{enumitem}
\usepackage{xcolor}
\usepackage{url}
\usepackage{titlesec}
\usepackage{multicol}
\setlength{\parskip}{0.28em}
\setlength{\parindent}{0pt}
\titlespacing*{\section}{0pt}{0.45em}{0.25em}
\titleformat{\section}{\bfseries\normalsize}{}{0em}{}
\pagestyle{empty}
\AtBeginDocument{\fontsize{8.5}{11.5}\selectfont}

\definecolor{accent}{HTML}{1a5276}

\begin{document}

\begin{center}
  {\Large\bfseries\color{accent} 单迹｜订单协同智能体}\\[0.15em]
  {\footnotesize 首届瓯客松 AI 黑客松（温州场）\;\textbar\;赛道：\textbf{第 07 题「碎片化沟通中的订单与任务协同」}}
\end{center}

\section*{目标用户}
以聊天、表格、电话推进业务的\textbf{中小企业}团队：销售、生产、订单负责人、成品仓与跟单。业务变化分散在沟通碎片中，需要把随口要求变成有人负责、能持续跟进的行动。

\section*{核心价值}
命题所强调的现实是：订单变化藏在聊天记录里，交付承诺出现在电话中，交接遗漏的一句话会让工作停住。许多中小企业并非靠完整系统推进，而是靠不断发生的沟通——消息很多，待跟进事项却易被淹没。

单迹从聊天、订单与零散表格（及可选通话转写、图片、录音）中\textbf{理解业务变化}，抽取结构化事件草案，对照订单识别数量/交期/库存冲突，生成带责任人、截止时间与证据的待办，并在新消息进入时\textbf{关联已有订单与未关闭任务}。模型只提建议；\textbf{写入订单与任务须人工确认}——让 AI 不只总结说了什么，而是帮助团队把事情做下去。

\section*{主要功能与亮点}
\begin{multicols}{2}
\begin{itemize}[leftmargin=1em,itemsep=0.08em,topsep=0.15em]
  \item \textbf{全局协作 AI}：跨工作台、审核、任务、仓库调度；改库前仍走确认卡片。
  \item \textbf{多模态接入}：聊天/表格/PDF/DOCX；可选 OCR、ASR；转写与聊天同链路解析。
  \item \textbf{变更识别与风控}：数量、交期、承诺、取消、线下洽谈；规则校验库存与产能冲突。
  \item \textbf{跨时间任务跟踪}：时间线联动；创建/更新/关闭建议；审计与修改前后记录。
  \item \textbf{智能库存管理}：共用在手、占用与可承诺量；多订单冲突提示；承诺不擅自改单。
  \item \textbf{严谨提醒与证据}：风险/任务一键回看源消息或单元格；采纳变更须显式确认。
  \item \textbf{全链路演示素材}：三阶段协作、缺客户订货表、文字与电话矛盾、沿用上次规格等。
\end{itemize}
\end{multicols}

\section*{作品访问与运行方式}
环境 Python 3.11+、Node.js。默认可无 API 密钥演示（fixture）；高级能力配置 \texttt{backend/.env}。

\noindent\textbf{后端：}\ \texttt{cd backend} $\rightarrow$ \texttt{pip install -r requirements.txt} $\rightarrow$ \texttt{uvicorn app.main:app --reload --port 8000}

\noindent\textbf{前端：}\ \texttt{cd frontend} $\rightarrow$ \texttt{npm install} $\rightarrow$ \texttt{npm run dev}

\noindent\textbf{访问：}\ \url{http://127.0.0.1:5173}。建议打开 \textbf{ORD-018 青禾成套}，粘贴 README 三段聊天并「解析为草案」，在审核页查看冲突、任务与证据；可继续体验任务看板、仓库与全局 AI 助手。

\end{document}
